\subsection{Ordinary relative coefficient deformation in classes two and three}

% Mathematical recovery transcription; see recovery provenance below.
\paragraph{Recovery provenance.}
This is a complete mathematical reconstruction of E.216, not a byte-identical recovery of its lost TeX. The final pre-outage proof was fully read positively by root, ghost\_audit, and literature. The two statement repairs required by those reviews are included: the prescribed higher B-cocycle's ramification belongs to $S$, and finite iteration requires compatible old local lifts through the requested precision. No mathematical computation or formal certificate was created for E.216. The pre-outage isolated TeX compilation was only a document check.

\subsubsection*{Notation and statement}
Let $N_r$ be the free geometric pro-five nilpotent group of class $r$, with the actual tangential arithmetic action of $G_\Q$ for the tripod. For $r=2,3$, integral Lazard coordinates identify its coefficient quotients with the free nilpotent Lie lattice modulo $5^n$. All arithmetic path coefficients and local Frobenius frames below are actual ones.

Fix $n\ge1$ and an actual global coefficient-$5^n$ class-$r$ section $s_n$. Let $a,b$ denote its residual abelian Kummer classes. In the full deformation problem the tangent module is the Lie algebra
\[
 M=\mathfrak n_r\otimes\F_5
\]
with action twisted by $s_n$. In the A-only problem prescribe a continuous global higher second-coordinate Kummer cocycle $B_{n+1}$, with its cocycle law, reducing literally to the second coordinate of $s_n$. The tangent module is then the ideal
\[
 I=I_X\otimes\F_5=\ker(M\longrightarrow\F_5Y).
\]
The word ``literal'' refers to the chosen representatives of the retained finite section and of its abelian cocycles; local geometric aligning frames remain available.

\begin{result}{E.216}{Proved}{Finite ordinary relative deformation in classes two and three}
For $r=2$ assume $(a,b)\ne(0,0)$ in the full problem, and assume $b\ne0$ in the A-only problem. For $r=3$ assume $a,b$ are independent over $\F_5$. Let $S$ contain $5,\infty$, all ramification of $s_n$, and, in the A-only problem, all ramification of the prescribed literal cocycle $B_{n+1}$.

Suppose next-coefficient local lifts are supplied at every place in $S$, retaining the chosen lower section and, in the A-only problem, the prescribed literal higher B-cocycle. Then a global next-coefficient lift exists matching these local classes. Its additional ramification can be confined to finitely many arbitrarily large primes $q$ with $\operatorname{ord}(q\bmod5)=4$. At each such prime the lift is the image of an actual ordinary point section: an A-point $q^m$ in the A-only problem, and an assigned A- or B-point $q^m$ or $1-q^m$ in the full problem. The integer $m$ may be chosen positive, divisible by $5^n$, and divisible by any specified integer prime to five.

There is also a retained-arithmetic version. A finite Galois field $F/\Q$ whose commutator subgroup has order prime to five may be retained, with a specified compatible arithmetic Frobenius $\sigma$ having primitive cyclotomic image modulo five. In this version the possible invisible arithmetic characters are explicitly retained. The interpolation conclusion holds when their reciprocity pairing with the prescribed abelian tangent targets is zero; it suffices that those targets be restrictions of a single global abelian tangent class, unramified outside the chosen old set. In the A-only problem this is a single first-coordinate compatibility condition.

Finite iteration is permitted when compatible old local lifts are supplied at every step up to the requested precision. In particular it is permitted for targets that are reductions of fixed integral local sections, including fixed ordinary points. An arbitrary finite old local lift is not asserted to extend indefinitely.
\end{result}

\begin{proof}
\par\medskip\noindent\textbf{1. The actual finite tangent modules}\par\nobreak\smallskip
Put
\[
 C=[X,Y],\qquad D=[C,X],\qquad E=[C,Y].
\]
In class three the degree-two arithmetic path coefficient is exactly
\[
 \eta(g)=\frac{\chi(g)^2-1}{24}.
 \tag{1}
\]
The fixed integral peripheral coordinate $Y^*=\operatorname{Ad}(\exp(C/24))Y=Y+E/24$ removes this coboundary contribution and is an eigen-coordinate for the order-four cyclotomic action. Equivalently one may keep (1) explicitly; the two descriptions are isomorphic actual modules. There is no additional degree-two arithmetic IA generator on the kernel of $\bar\chi$. Degree-three path coefficients act beyond class three and are invisible here.

In class two, the geometric action on the tangent module is the translation image $V\le\F_5(1)^2$ supplied by $(a,b)$. In the full problem $V\ne0$; in the A-only problem its B-projection is nonzero. In class three, independence of $a,b$ gives the full class-two geometric image by the Frattini criterion. The actual minimal action group is thus the appropriate $V\rtimes C_4$ or $P_2\rtimes C_4$, where $P_2$ is the free class-two exponent-five group. These are statements about the image of the retained section, not assumptions that $G_\Q$ is projective.

Let $D_M=M^*(1)$ or $D_I=I^*(1)$ be the Tate dual. Their geometric filtrations have constant graded pieces of cyclotomic weights
\[
 0,-1 \quad(r=2),\qquad 0,-1,-2 \quad(r=3),
 \tag{2}
\]
with the evident multiplicities. The geometric abelianization has weight one. Hence for each graded piece $W$ in (2),
\[
 H^1(P,W)^{C_4}=\operatorname{Hom}(P^{\mathrm{ab}},W)^{C_4}=0.
\]
The order of $C_4$ is prime to five, so its positive cohomology vanishes. Induction through the filtration therefore proves
\[
 H^1(J,D_M)=0\quad\text{or}\quad H^1(J,D_I)=0
 \tag{3}
\]
for the actual minimal action group $J$. No rank-two assumption is needed for the class-two full-module calculation.

\par\medskip\noindent\textbf{2. An unrestricted global reference lift}\par\nobreak\smallskip
The next-coefficient extension has the indicated finite abelian tangent kernel. Its obstruction is a class in $H^2(G_{\Q,S},M)$ or $H^2(G_{\Q,S},I)$. It is zero locally at the old places by the supplied local lifts. At a place outside $S$ it has an unramified local lift: lift the finite Frobenius element and send inertia to one. In the fixed-B problem this assertion uses explicitly that $B_{n+1}$ is unramified outside $S$.

We justify passage to a global weak solution without assuming an unrestricted embedding theorem has ordinary ramification. By Chebotarev, a global finite dual class which is zero at every decomposition group descends to the finite action group: on the kernel of the action its restriction is a homomorphism, and Frobenius detection kills it. Equation (3) then kills the descended class. Equivalently the relevant cyclic-restriction kernel of $J$ is zero.

Choose finitely many temporary places $T_0$ detecting the finite vector space of possible dual classes over $G_{\Q,S}$. A class over $G_{\Q,S\cup T_0}$ that is locally zero at $T_0$ has zero new inertia and descends to $G_{\Q,S}$; detection then kills it. Finite Poitou--Tate duality gives the corresponding vanishing of the degree-two local-global kernel. The obstruction therefore vanishes after allowing $T_0$, and a global weak next lift exists. We later prescribe unramified targets at $T_0$ and remove its unwanted temporary ramification.

This argument is compatible with the original Neukirch localization criterion: for a finite module $A$ with dual $D$, his Corollary 6.4(c) gives finite-set localization surjectivity when $H^1(\operatorname{Gal}(\Q(D)/\Q),D)=0$. Here we have computed that hypothesis for the actual module, rather than assuming it from solvability of the kernel.

\par\medskip\noindent\textbf{3. Ordinary local tangent lines and Frobenius}\par\nobreak\smallskip
At an auxiliary prime $q$ with primitive $\bar\chi(F_q)$, the retained lower section is unramified. Since
\[
 q^i-1\in\Z_5^\times\qquad(1\le i\le r\le3),
\]
successive central lifting shows that every unramified lower class is geometrically conjugate to the actual cusp graph, with no residual unramified tangent ambiguity. In particular this is an actual local statement including Frobenius.

The local tangent cohomology is its tame residue space:
\[
 H^1(\Q_q,M)=M(-1)^{F_q},\qquad
 H^1(\Q_q,I)=I(-1)^{F_q}.
 \tag{4}
\]
Only the weight-one directions survive in (4). For the full module they are the A- and B-lines; for the ideal there is only the A-line. A point $q^m$ with $m\equiv5^n a\pmod{5^{n+1}}$ realizes the A-line, and $1-q^m$ realizes the B-line. Because unramified tangent cohomology is zero, equality of the tame class implies equality of the complete local section class, including Frobenius. The local class depends only on this integer valuation at such a prime; pro-five unit Kummer classes vanish since $q-1$ is a five-adic unit.

In the fixed-B problem the prescribed higher B-cocycle has zero local Kummer class. If its chosen cochain is not literally zero, first adjust the representative of the ordinary point by a congruence Y-alignment. This retains the lower frame and makes the higher B-cocycle literally prescribed. The remaining comparison is then exactly the $I$-torsor.

The dual local cohomology has no tame residue, since
\[
 D_M(-1)^{F_q}=0,\qquad D_I(-1)^{F_q}=0,
 \tag{5}
\]
all its graded weights being $-1,-2,-3$. Thus allowing such auxiliary primes creates no new dual classes with admissible local conditions. This is the crucial low-class fact; it is not asserted in all weights.

\par\medskip\noindent\textbf{4. Detection by ordinary auxiliary directions}\par\nobreak\smallskip
Fix a dual class over the old finite ramification set. Consider the finite extension carrying the actual action, the retained section, and that dual cocycle. If its local pairing vanishes at every permitted ordinary A- or B-prime, compare Frobenius elements $g$ and $gn$ with $n$ in the kernel of the retained finite representation. The local Tate edge pairing says that the value on $n$ annihilates the corresponding conjugated peripheral tangent line.

These conjugated lines span the required module. In class two, a nonzero translation image and the two peripheral directions generate $X,Y,C$. In the A-only class-two problem the nonzero B-translation generates $C$ from $X$. In class three, the full class-two geometric image gives the whole free tangent module, or the ideal generated by $X$, by iterated adjoint brackets. The finite exponential coefficients involved have denominators only $2$ and $3$, which are units modulo five. Consequently the dual class vanishes on the kernel and descends to the minimal finite action group. Equation (3) kills it.

There are only finitely many old dual classes. Hence finitely many permitted auxiliary primes suffice to kill them all; Chebotarev permits these primes to avoid any prescribed finite set and exceed any prescribed bound. At each auxiliary prime select the appropriate ordinary A- or B-line. By (5), any dual class with the resulting local conditions over the enlarged ramification set descends to the old set, so no new undetected dual class appears.

Finite Poitou--Tate duality now identifies the image of global tangent corrections with the annihilator of the dual obstruction. The latter is zero. We can therefore realize the prescribed old tangent discrepancies while retaining the selected ordinary lines at all auxiliary primes. Twisting the global reference lift gives the desired section. Include the temporary set $T_0$ among the old places with unramified targets; this removes the temporary ramification introduced in the existence argument.

\par\medskip\noindent\textbf{5. Retained arithmetic fields and the exact remaining compatibility}\par\nobreak\smallskip
Suppose a finite field $F$ with commutator of order prime to five is retained. Enlarge it by the required cyclotomic field if necessary. Its intersection with the residual geometric-action extension is arithmetic: the latter's geometric five-group lies in the derived subgroup, so it cannot map nontrivially into the commutator of $\operatorname{Gal}(F/\Q)$. Therefore full geometric image survives above each compatible arithmetic element. The prescribed Frobenius $\sigma$ must first be an actual compatible element of the whole retained compositum. A group-order assertion alone would not establish compatibility with independent radical-root conditions.

Repeat the preceding kernel comparison using only Frobenii above $\sigma$. The same peripheral spans kill the geometric kernel. Inflation--restriction and (3) leave exactly arithmetic characters with values in the constant abelian dual:
\[
 \psi\in\operatorname{Hom}(\operatorname{Gal}(F'/\Q),\F_5)^d,
 \qquad \psi(\sigma)=0,
 \tag{6}
\]
where $d=2$ for the full module and $d=1$ for the fixed-B ideal; $F'$ denotes the actual retained arithmetic compositum. We do not silently discard (6).

Their pairing with an old tangent discrepancy is the usual arithmetic cup pairing with its abelian projection. If that projection is the localization of one global abelian tangent class, unramified outside the old set, global reciprocity makes the sum of the old local pairings zero. At the auxiliary places (6) gives zero, and outside the chosen set both classes are unramified. Thus the old targets annihilate the exact invisible arithmetic kernel, and Poitou--Tate again supplies the required correction.

One may start this retained-field argument from the unrestricted temporary reference lift. At its additional temporary places choose compatible unramified targets with the same global abelian tangent class, then remove those places by the same correction argument. This preserves the stated compatibility rather than assuming arbitrary targets are attainable under arbitrary extra arithmetic restrictions.

\par\medskip\noindent\textbf{6. Integer representatives and finite iteration}\par\nobreak\smallskip
At a new place the residue fixes $m$ only modulo $5^{n+1}$, with $5^n\mid m$. Choose a positive integer representative also divisible by any prescribed prime-to-five integer; the Chinese remainder theorem gives such representatives, arbitrarily large if needed. Freeze that integer and retain the corresponding actual ordinary local point at subsequent coefficient steps.

For a finite sequence of coefficient refinements, retain every previously introduced ordinary point as an old local target. These new targets lift integrally by construction. The original old targets must likewise have supplied compatible local lifts through the desired precision; this holds if they came from fixed integral local sections. At the coefficient place five, when fixing the abelian data in classes two and three, the successive local central obstructions vanish because $H^2(\Q_5,\Z/5^k(w))=0$ for $w=2,3$ by local duality. The real positive cohomology terms vanish at this odd coefficient.

This proves finite iteration under the stated local hypotheses. It does not change a completed ordinary pair by finitely many rational multipliers, establish an ordinary infinite tail, or supply the initial lift from one nilpotent class to the next. Those are separate problems.
\end{proof}

\subsubsection*{Primary source locators and limits}
The local pairing and finite global duality used above are Milne, \emph{Arithmetic Duality Theorems}, Chapter I, Sections 2 and 4, especially Theorems 4.9--4.10, pp.~56--57. The original finite-module localization criterion was independently opened in Neukirch (1973), Corollary 6.4(c), printed p.~90, in the complete volume available at
\url{https://gdz.sub.uni-goettingen.de/download/pdf/PPN356556735_0021/PPN356556735_0021.pdf}.
Its arbitrary-auxiliary-ramification conclusion alone is insufficient for the ordinary-prime assertion; Sections 3--5 above are the additional module and local-condition proof. No projectivity of $G_\Q$, automatic class jump, or infinite compatible branch is asserted.
